\title{FlashAttention-3: Fast and Accurate Attention with Asynchrony and Low-precision}

\begin{abstract}
Attention mechanisms form the foundational layer of the widely adopted Transformer model, and often represent the principal computational bottleneck for scaling large language models and extending sequence length. The \textsc{FlashAttention} algorithm introduced optimized attention computation on GPUs by reducing redundant memory operations. Nonetheless, prior versions did not fully harness advancements in modern GPU hardware: for instance, \textsc{FlashAttention-2} makes use of only 35\% of available resources on the H100 GPU. In this work, we present three key innovations to enhance attention performance on NVIDIA Hopper GPUs: leveraging the asynchronous capabilities of Tensor Cores and TMA, we (1) enable simultaneous data transfer and computation through warp-specialization, (2) combine block-level matrix multiplication and softmax steps in an interleaved fashion, and (3) implement block quantization and non-coherent computation to capitalize on FP8 low-precision support. Empirically, our new approach, \textsc{FlashAttention-3}, realizes a 1.5-2.0$\times$ acceleration on H100 GPUs, attaining up to 740 TFLOPs/s in FP16 (75\% hardware utilization) and nearly 1.2 PFLOPs/s in FP8. Furthermore, we find that FP8 \textsc{FlashAttention-3} produces numerical errors 2.6$\times$ lower than a standard FP8-based attention baseline.
\end{abstract}

\section{Introduction}
\label{sec:intro}

Within the Transformer architecture~\citep{vaswani2017attention}, the principal computational challenge arises from the attention mechanism, as the process of calculating self-attention scores between queries and keys exhibits quadratic complexity with respect to the input sequence length. Enhancing the scalability of attention mechanisms to support extended contexts could enable novel functionalities, such as modeling and reasoning across multiple lengthy documents~\citep{guo2021longt5,shaham2022scrolls,peng2023yarn}, handling extensive code repositories~\citep{roziere2023code, li2023starcoder}, processing new data modalities (e.g., high-resolution visual inputs~\citep{chen2022scaling}, auditory signals~\citep{gulati2020conformer}, and video~\citep{ho2022video}), as well as addressing emerging tasks (such as modeling user interactions involving extensive histories~\citep{sun2019bert4rec} or facilitating agent workflows across long horizons~\citep{yao2022react}). These considerations have sparked widespread efforts aimed at improving attention efficiency for long-context applications, including approaches based on approximate computation~\citep{katharopoulos2020transformers,choromanski2020rethinking,
tay2020efficient}, advances in software implementations (\citep{rabe2021self,dao2022flashattention,kwon2023efficient}), and proposals for entirely different neural architectures~\citep{peng2023rwkv,sun2023retentive,gu2023mamba}.

This paper extends the contributions of \citet{dao2022flashattention}, who developed exact-attention algorithms by embedding an understanding of GPU execution patterns and hardware properties at the algorithmic design stage. Dao et al., in \citep{dao2022flashattention}, presented \textsc{FlashAttention}, which employed an innovative tiling approach to parallelize the attention mechanism. Their approach avoids redundant accesses to slower global memory by consolidating all attention-related operations within a unified GPU kernel.

Subsequently, \citet{dao2023flashattention2} refined the method by introducing \textsc{FlashAttention-2}, which further parallelizes computations across the sequence length axis and processes the core loop in blocks of the key and value matrices. This modification enhances resource allocation and task distribution on the GPU. Nevertheless, our analysis reveals that, when deployed on modern GPUs, \textsc{FlashAttention-2} still underutilizes available computational capacity compared to highly optimized matrix-multiplication (GEMM) routines—achieving, for instance, only 35\% as opposed to the 80-90\% observed with GEMM on the Hopper H100 hardware. Some of this discrepancy can be traced to technical implementation issues, such as a failure to employ instructions specific to the Hopper architecture and instead defaulting to those intended for the previous Ampere architecture when leveraging Tensor Cores. Recent efforts—including ThunkerKitten~\citep{spector2024thunder} and cuDNN 9~\citep{cudnn9}—have demonstrated that using Hopper-optimized instructions alongside tile-centric abstractions not only accelerates attention workloads but also streamlines the associated codebase.

At a more foundational level, the algorithm behind \textsc{FlashAttention-2} operates within a streamlined synchronous paradigm, deliberately omitting explicit mechanisms for asynchrony and low-precision computation. Asynchrony, however, naturally emerges from contemporary hardware architectures that target the acceleration of core machine learning tasks: this includes specialized units for matrix multiplication (such as Tensor Cores) and for memory access (Tensor Memory Accelerator -- TMA), which work in parallel with general-purpose CUDA cores responsible for logic, integer, and standard floating-point operations. The adoption of low-precision formats, for instance FP8 in Hopper and FP4 in Blackwell—building on the precedent set by FP16 (Pascal, 2017) and BF16 (Ampere, 2020)—has consistently yielded significantly higher throughput (by factors of two or four) without increasing power consumption or chip real estate. We examine these architectural advances in Hopper and their relevance in \cref{subsec:hardware}.

To address these developments, the core technical hurdle lies in reengineering \textsc{FlashAttention-2} so it can leverage such hardware optimizations: specifically, asynchrony introduces the need to interleave matmul and softmax operations, despite their inherent data dependencies, while the move to lower precisions demands rigorous strategies to control quantization error, a concern made more acute by the occurrence of outlier features in large language models~\citep{dettmers2208llm, sun2024massive}.

To address this, we introduce \textsc{FlashAttention-3}, which integrates three novel strategies designed to enhance performance on modern GPU platforms:\footnote{Our findings are presented with respect to NVIDIA's Hopper architecture, though the proposed algorithm remains applicable to any GPU with advanced asynchronous execution and support for low-precision operations.}

\begin{enumerate}

\item \textbf{Producer-Consumer asynchrony:} We introduce a warp-specialized software pipelining method that takes advantage of asynchronous data transfers and Tensor Core operations. By allocating data producer and consumer roles to distinct warps, our approach broadens the potential to mask both memory access and instruction issuance delays within the algorithm.

\item \textbf{Hiding softmax under asynchronous block-wise GEMMs:} Our design overlaps the execution of relatively low-throughput operations necessary for softmax, including floating-point multiply-adds and exponential functions, with asynchronous WGMMA instructions performing GEMM. To facilitate this overlap, we restructure the \textsc{FlashAttention-2} algorithm, eliminating some of the sequential dependencies between softmax computation and GEMM execution. Specifically, in the 2-stage variant of our scheme, one block of the score matrix is being processed by softmax while the WGMMA unit asynchronously prepares the subsequent block.

\item \textbf{Hardware-accelerated low-precision GEMM:} We modify the forward pass routine to target GEMM execution on FP8 Tensor Cores, which nearly doubles empirical TFLOPs/s rates. Achieving this performance necessitates reconciling the differing memory layout expectations imposed by WGMMA regarding FP32 accumulators and FP8 operand matrices. To address the reduction in precision and attendant accuracy loss from using FP8, we apply block quantization along with incoherent computation strategies.

To empirically assess the effectiveness of our approach, we conduct a series of benchmarks for \textsc{FlashAttention-3} using the H100 SXM5 GPU, exploring various parameter settings. Our results demonstrate that (1) the FP16 mode provides a 1.5-2.0$\times$ acceleration in the forward pass—achieving up to 740 TFLOPs/s—and a 1.5-1.75$\times$ improvement in the backward pass, relative to \textsc{FlashAttention-2}, (2) FP8 mode attains near 1.2 PFLOPs/s, and (3) at greater sequence lengths, FP16 delivers superior performance while FP8 shows strong competitiveness\footnote{Specifically, FP8 in \textsc{FlashAttention-3} leads for a head dimension of 64, and for dimensions 128 and 256, it matches performance where causal masking is absent, while trailing slightly with causal masking enabled.} against NVIDIA cuDNN's state-of-the-art attention implementation. Further, we verify that \textsc{FlashAttention-3} in FP16 incurs the same level of numerical error as \textsc{FlashAttention-2}, outperforming standard attention in accuracy due to retention of intermediate values (such as those from softmax rescaling) in FP32 precision. Additionally, in scenarios characterized by outlier features, the FP8 variant of \textsc{FlashAttention-3} with block quantization and incoherent processing yields accuracy that is 2.6$\times$ higher than standard attention that employs per-tensor quantization.

\textsc{FlashAttention-3} has been released under a permissive open-source license,\footnote{\textsc{FlashAttention-3} can be found at \texttt{https://github.com/Dao-AILab/flash-attention}} and there are ongoing efforts to incorporate it into PyTorch and Hugging Face libraries, aiming to maximize accessibility for the research and developer community.

\section{Background: Multi-Head Attention and GPU Characteristics}
\label{sec:background}

\subsection{Multi-Head Attention}
\label{subsec:multi_head_attn}

Let $\mathbf{Q}, \mathbf{K}, \mathbf{V} \in \mathbb{R}^{N \times d}$ be the query, key and value input sequences associated to a single head, where $N$ is the sequence length and $d$ is the head dimension. Then the attention output $\mathbf{O}$ is computed as:

\begin{equation*}
  \mathbf{S} = \alpha \mathbf{Q} \mathbf{K}^\top \in \mathbb{R}^{N \times N}, \quad \mathbf{P} = \mathrm{softmax}(\mathbf{S}) \in \mathbb{R}^{N \times N}, \quad \mathbf{O} = \mathbf{P}\mathbf{V} \in \mathbb{R}^{N \times d},
\end{equation*}

Here, $\mathrm{softmax}$ is executed along each row, and the scaling parameter $\alpha$ is commonly chosen as $1/\sqrt{d}$. To maintain numerical stability when working with exponentials, it is standard practice to subtract $\mathrm{rowmax}(\mathbf{S})$ from $\mathbf{S}$ before applying the exponential. In the context of multi-head attention (MHA), individual heads use distinct projections for queries, keys, and values. This process is computed in parallel across all heads and batches, yielding the final output tensor.

Now let $\phi$ be a scalar loss function and let $\mathbf{d}(-) = \partial \phi / \partial (-)$ be notation for the gradient. Given the output gradient $\mathbf{dO} \in \mathbb{R}^{N \times d}$, we compute $\mathbf{dQ}$, $\mathbf{dK}$, and $\mathbf{dV}$ according to the chain rule as follows:

\begin{align*}
  \mathbf{dV} &= \mathbf{P}^\top \mathbf{dO} \in \mathbb{R}^{N \times d} \\
  \mathbf{dP} &= \mathbf{dO} \mathbf{V}^\top \in \mathbb{R}^{N \times N} \\
  \mathbf{dS} &= \mathrm{dsoftmax} (\mathbf{dP}) \in \mathbb{R}^{N \times N} \\
  \mathbf{dQ} &= \alpha \mathbf{dS} \mathbf{K} \in \mathbb{R}^{N \times d} \\
  \mathbf{dK} &= \alpha \mathbf{dS}^\top \mathbf{Q} \in \mathbb{R}^{N \times d},
\end{align*}

In this context, $\mathbf{d}s = (\mathrm{diag}(p) - p p^\top)\mathbf{d}p$, where $p = \mathrm{softmax}(s)$ with respect to the vector $s$, and the notation $\mathrm{dsoftmax}(\mathbf{dP})$ denotes applying this operation to each row independently. This calculation is also efficiently parallelized over multiple heads and batches during the backward phase of MHA.

\subsection{GPU hardware characteristics and execution model}
\label{subsec:hardware}

We outline the elements of the GPU execution model pertinent to \textsc{FlashAttention-3}, centering our discussion on the NVIDIA Hopper architecture as a specific example of this framework.

\paragraph{Memory hierarchy:} The GPU architecture features a layered memory system, where the amount of storage typically decreases as memory access speed increases (\cref{tab:gpu-hierarchy})\footnote{\citet{luo2024benchmarking} gives a figure of 128 bytes per clock cycle in shared memory bandwidth for each SM; multiplying this by 132 SMs and a boost clock frequency of 1830 MHz yields the aggregate bandwidth.}. The most distant layer is global memory (GMEM), also referred to as HBM, which consists of off-chip DRAM shared by all streaming multiprocessors (SMs). GMEM is automatically backed by an on-chip L2 cache, which improves data locality. Within each SM, there is a modest-sized, on-chip cache known as shared memory (SMEM), which is managed directly by the programmer and organized into multiple banks for parallel access. At the innermost level lies the register file, unique to each SM.

\paragraph{Thread hierarchy:} In the GPU programming paradigm, execution units known as threads are structured into several hierarchical layers. Starting from the most granular level, threads are grouped into warps (each containing 32 threads), which are then organized into warpgroups (comprising 4 adjacent warps). These are further aggregated into threadblocks (also referred to as cooperative thread arrays or CTAs), then into threadblock clusters (available in Hopper architectures), and finally into grids at the highest level.

The two hierarchies exhibit a strong interconnection. Threads belonging to a particular CTA are executed together on a single SM, while CTAs assigned to the same cluster are simultaneously scheduled on one GPC. All threads within a CTA have direct access to SMEM, whereas each individual thread possesses up to 256 private registers (RMEM), inaccessible to other threads.

\paragraph{Asynchrony and warp-specialization:}

GPUs function as throughput-oriented processors, leveraging parallelism and asynchronous operations to mitigate both memory and execution delays. To support asynchronous data transfers between GMEM and SMEM, the Hopper architecture introduces the Tensor Memory Accelerator (TMA), a specialized hardware component \cite[\S7.29]{cuda}. Moreover, distinguishing itself from earlier architectures like Ampere, Hopper’s Tensor Core—accessible through the warpgroup-wide WGMMA instruction \cite[\S9.7.14]{ptx}—operates asynchronously as well and is capable of receiving input data directly from shared memory.

Enabling hardware-level asynchrony facilitates the implementation of warp-specialized kernels, where the warps within a CTA are explicitly assigned as either producers or consumers, restricting their operations to data transfers or computational tasks, respectively. This architectural distinction generally enhances the compiler's effectiveness in producing efficient instruction schedules \citep{warp-specialization-2011}. Furthermore, the Hopper architecture introduces the capability to dynamically redistribute registers among different warpgroups using the \verb|setmaxnreg| command \cite[\S9.7.17.1]{ptx}, thus permitting warps engaged in MMAs to access a greater portion of RMEM compared to warps handling TMA, which typically require just a single thread.

\paragraph{Low-precision number formats:}
\label{sec:low-precision-gpu}
Recent GPUs incorporate dedicated hardware to enhance performance for computations using low-precision formats. As an illustration, on Hopper architecture, the FP8 Tensor Cores can be utilized via the WGMMA instruction, resulting in double the throughput per SM relative to computations in FP16 or BF16.

Correct use of FP8 WGMMA requires familiarity with the memory layout requirements imposed on the input matrices. For instance, consider a GEMM operation computing $A \times B^{\top}$ where $A$ is of size $M \times K$ and $B$ is $N \times K$. If either $A$ or $B$ has elements arranged contiguously along the $M$ or $N$ axis, it is termed \emph{mn-major}; conversely, it is described as \emph{k-major} when elements are contiguous in the $K$ dimension. While FP16 WGMMA permits both mn-major and k-major memory layouts for input operands stored in SMEM, the FP8 WGMMA restricts allowable formats exclusively to the k-major layout. This limitation can present challenges in applications—such as attention mechanisms—where sequential GEMMs are fused within a single kernel, since mismatched layouts between FP32 accumulators and FP8 input operands can prevent successful chaining of dependent FP8 WGMMA calls.

Within the realm of attention mechanisms, such layout constraints necessitate specific adjustments to the architecture of an FP8 algorithm; these adaptations are detailed in \cref{sec:algofp8}.

\subsection{Standard Attention and Flash Attention} In line with the description provided by \citet{dao2022flashattention}, we define \textbf{standard attention} as an attention operation on the GPU that stores the intermediate results $\mathbf{S}$ and $\mathbf{P}$ in high-bandwidth memory (HBM). \textsc{FlashAttention} introduces a novel approach by employing a localized version of the softmax reduction, thus eliminating the costly process of reading and writing these intermediate matrices, and instead merging the attention computation into a single GPU kernel. This localized softmax process is reflected in lines \ref{code-ws:softmax_start}-\ref{code-ws:softmax_end} within the consumer mainloop of \cref{alg:flash3_wgmma_ws_only}, and also includes the rescaling operations applied to submatrices of $\mathbf{O}$. A concise proof showing that this strategy correctly computes $\mathbf{O}$ is available in \cite[\S 2.3.1]{dao2023flashattention2}.

\section{FlashAttention-3: Algorithm}
\label{sec:algo}

This section presents an overview of the \textsc{FlashAttention-3} algorithm. To maintain clarity, our primary emphasis is on the forward computation; the procedure for the backward pass is detailed in~\cref{sec:algo_ws_bwd}. Initially, we discuss the incorporation of warp-specialization combined with a circular SMEM buffer into the foundational \textsc{FlashAttention-2} method. Subsequently, we outline how leveraging the asynchrony characteristic of WGMMA enables us to construct a two-stage GEMM-softmax pipeline with overlapping computation. In the final part of the section, we detail the adjustments necessary for supporting FP8, addressing both layout compatibility and enhancements to numerical precision through the use of block quantization and incoherent computation strategies.

\subsection{Producer-Consumer asynchrony through warp-specialization and pingpong scheduling}
\label{sec:algo_ws}

\paragraph{Warp-specialization}
Similar to \textsc{FlashAttention-2}, the forward propagation in \textsc{FlashAttention-3} exhibits extensive parallelism across batch size, the number of heads, and the length of the query sequence. Accordingly, we focus on the algorithm from a CTA-level perspective, in which a specific tile $\mathbf{Q}_i$ of the query matrix is mapped to compute the corresponding output tile $\mathbf{O}_i$. For clarity, the following exposition describes the warp-specialization method using a circular SMEM buffer, while temporarily omitting the mechanism for overlapping GEMM with the softmax step. Denoting the head dimension by $d$ and the sequence length by $N$, we establish a query block size $B_r$ to partition $\mathbf{Q}$ into $T_r = \lceil \frac{N}{B_r} \rceil$ blocks, namely $\mathbf{Q}_1, .., \mathbf{Q}_{T_r}$.

\begin{algorithm}[H]
    \caption{\small\label{alg:flash3_wgmma_ws_only}\textsc{FlashAttention-3} forward pass \textbf{excluding} intra-consumer overlap -- CTA view}
    \begin{algorithmic}[1]
\REQUIRE Matrices $\mathbf{Q}_i \in \mathbb{R}^{B_r \times d}$ and $\mathbf{K}, \mathbf{V} \in \mathbb{R}^{N \times d}$ stored in HBM, key block size $B_c$ such that $T_c = \lceil \frac{N}{B_c} \rceil$.
\STATE Create a pipeline manager to coordinate barriers using a circular buffer in shared memory of $s$ stages.
\IF {executing as the producer warpgroup}
\STATE Release a specified count of registers.
\STATE Load $\mathbf{Q}_i$ from HBM into shared memory.
\STATE Once done, signal the consumer that $\mathbf{Q}_i$ load is complete.
\FOR{$0 \le j < T_c$}
    \STATE Wait for consumption of stage $(j\,\%\,s)$ of the buffer.
    \STATE Load $\mathbf{K}_j$ and $\mathbf{V}_j$ from HBM to the corresponding buffer stage $(j\,\%\,s)$ in shared memory.
    \STATE On load completion, notify the consumer of $\mathbf{K}_j, \mathbf{V}_j$ readiness.
\ENDFOR
\ELSE \STATE Allocate a computed number of registers based on the consumer warp count.
\STATE Initialize on-chip: $\mathbf{O}_i = (0) \in \mathbb{R}^{B_r \times d}$; $\ell_i = 0$, $m_i = -\infty$, both in $\mathbb{R}^{B_r}$.
\STATE Await availability of $\mathbf{Q}_i$ in shared memory.
\FOR{$0 \le j < T_c$}
\STATE Wait until $\mathbf{K}_j$ becomes available in shared memory.
\STATE Perform matrix multiplication $\mathbf{S}_i^{(j)} = \mathbf{Q}_i \mathbf{K}_j^T$ (using SS-GEMM). Commit, then await completion. \label{alg:ws_only_gemm1}
\STATE Cache $m_i^{\mathrm{old}} = m_i$; update $m_i = \mathrm{max}(m_i^{\mathrm{old}}, \mathrm{rowmax}(\mathbf{S}_i^{(j)}))$. \label{code-ws:softmax_start}
\STATE Calculate $\widetilde{\mathbf{P}}_i^{(j)} = \mathrm{exp}(\mathbf{S}_i^{(j)} - m_i)$ and update $\ell_i = \mathrm{exp}(m_i^{\mathrm{old}} - m_i) \ell_i + \mathrm{rowsum}(\widetilde{\mathbf{P}}_i^{(j)})$. \label{code-ws:softmax_end}
\STATE Wait for $\mathbf{V}_j$ to become available in shared memory.
\STATE Compute $\mathbf{O}_i = \mathrm{diag}(\mathrm{exp}(m_i^{\mathrm{old}} - m_i))^{-1} \mathbf{O}_i + \widetilde{\mathbf{P}}_i^{(j)} \mathbf{V}_j$ (via RS-GEMM). Commit and synchronize. \label{alg:ws_only_gemm2}
\STATE Signal the producer that buffer stage $(j\,\%\,s)$ has been released.
\ENDFOR
\STATE Finalize by computing $\mathbf{O}_i = \mathrm{diag}(\ell_i)^{-1} \mathbf{O}_i$, and $L_i = m_i + \log(\ell_i)$.
\STATE Output $\mathbf{O}_i$ and $L_i$ to HBM as block $i$ of $\mathbf{O}$ and $L$ respectively.
\ENDIF
\end{algorithmic}
\end{algorithm}

In our deployment of \cref{alg:flash3_wgmma_ws_only} on the Hopper architecture, we utilize \verb|setmaxnreg| to manage (de)allocation of resources, employ TMA instructions for transferring $\mathbf{Q}_i$ as well as $\{ \mathbf{K}_j, \mathbf{V}_j \}_{0 \leq j < T_c}$, and leverage WGMMA for conducting GEMMs within the consumer mainloop; the SS or RS prefixes specify whether the primary input operand resides in shared memory or in the register file, respectively. When analyzing the operation of \cref{alg:flash3_wgmma_ws_only}, it is important to observe that TMA load commands do not block for the completion of preceding loads because they are handled asynchronously. Additionally, within the producer mainloop, buffer-filling ensures that no synchronization waits are enacted during the initial $s$ iterations.

\paragraph{Pingpong scheduling} The asynchrony provided by WGMMA and TMA, combined with warp specialization, enables overlapping the softmax computation occurring in one warpgroup with the GEMM execution in a different warpgroup. This strategy becomes appealing when considering that contemporary hardware accelerators deliver substantially higher throughput for matrix multiplication operations relative to non-matmul tasks. For instance, the H100 SXM5 GPU achieves 989 TFLOPS of FP16 matrix multiplication yet can only manage 3.9 TFLOPS for special function computations such as exponential\footnote{According to the CUDA programming guide, each streaming multiprocessor (SM) can carry out 16 special function operations per clock cycle. This yields 3.9 TFLOPS when multiplying 16 operations by 132 SMs and a clock speed of 1830 MHz.}. In the context of the FP16 attention forward pass with a head size of 128, the required number of matmul floating-point operations outpaces exponentials by a factor of 512, while the exponential throughput is 256 times slower, causing the exponential computation to consume up to 50\% of the overall execution time for those cycles. The discrepancy grows when using FP8, as matmul performance doubles but the throughput for exponentials remains unchanged.

Because the exponential operation is handled by a dedicated hardware component (the multi-function unit), it is most efficient to initiate the exponential computation concurrently with the matmul operations being executed on the Tensor Cores. To orchestrate this overlap, we introduce synchronization barriers (\texttt{bar.sync} instructions) that arrange for the GEMMs (specifically, GEMM1—$\mathbf{P} \mathbf{V}$ for a given iteration, and GEMM0—$\mathbf{Q} \mathbf{K}^\top$ for the subsequent iteration) assigned to warpgroup 1 to be launched prior to those assigned to warpgroup 2. Consequently, while warpgroup 2 proceeds with its GEMMs, the softmax for warpgroup 1 is being computed. After this, the warpgroups swap responsibilities, allowing warpgroup 2 to handle the softmax as warpgroup 1 returns to GEMMs, thus following a "pingpong" scheduling approach. This scheduling method is depicted in~\cref{fig:pingpong_scheduling}. Although in real workloads the alternation between warpgroups is not as perfectly synchronized as the diagram suggests, this strategy typically yields a noticeable performance enhancement—for instance, raising throughput from 570 TFLOPS to 620–640 TFLOPS on FP16 forward passes with a head dimension of 128 and sequence length of 8192.

\paragraph{Attention variants} In the cases of multi-query attention~\citep{shazeer2019fast} and grouped query attention~\citep{ainslie2023gqa}, we adopt the methodology outlined in \textsc{FlashAttention-2}, modifying tensor indexing so that $\mathbf{K}$ and $\mathbf{V}$ are not redundantly stored in HBM.

\subsection{Intra-warpgroup overlapping GEMMs and softmax}

It is possible to interleave certain instructions from the softmax operation with those from the GEMMs, even when working within a single warpgroup. Here, we outline a specific method to achieve this overlap.

In the attention algorithm, operations within the inner loop (main loop) have sequential dependencies that impede parallelization within a single iteration. For example, (local) softmax (lines \ref{code-ws:softmax_start} to \ref{code-ws:softmax_end}) relies on the output $\mathbf{S}_i^{(j)}$ of the first GEMM, while the second GEMM takes its result $\widetilde{\mathbf{P}}_i^{(j)}$ as an operand. Indeed, the wait statements in lines \ref{alg:ws_only_gemm1} and \ref{alg:ws_only_gemm2} of \cref{alg:flash3_wgmma_ws_only} serialize the execution of softmax and GEMMs. However, we can break these dependencies by pipelining across iterations through additional buffers in registers. Pursuing this idea, we propose the following two-stage\footnote{Note that the number of stages of the overlapping scheme is bounded by, but need not equal, the number $s$ of stages in the circular SMEM buffer.} GEMM-softmax pipelining algorithm:

\begin{algorithm}[H]
    \caption{\small\label{alg:flash3_wgmma}\textsc{FlashAttention-3} consumer warpgroup forward pass}
    \begin{algorithmic}[1]
      \REQUIRE  The matrices $\mathbf{Q}_i \in \mathbb{R}^{B_r \times d}$ and $\mathbf{K}, \mathbf{V} \in \mathbb{R}^{N \times d}$ are resident in HBM; define the key block size as $B_c$ and let $T_c = \lceil \frac{N}{B_c} \rceil$.
      \STATE Adjust allocation of registers dynamically, contingent on the quantity of consumer warps involved.
      \STATE Initialize on-chip data: set $\mathbf{O}_i = (0) \in \mathbb{R}^{B_r \times d}$, as well as $\ell_i, m_i = (0), (-\infty) \in \mathbb{R}^{B_r}$.
      \STATE Ensure both $\mathbf{Q}_i$ and $\mathbf{K}_0$ are fully loaded into shared memory before proceeding.
      \STATE Execute WGMMA to calculate $\mathbf{S}_{\mathrm{cur}} = \mathbf{Q}_i \mathbf{K}_0^T$. Commit operation and synchronize.
      \STATE Deallocate the $0$th segment of the buffer associated with $\mathbf{K}$.
      \STATE Determine $m_{i}$, $\tilde{\mathbf{P}}_{\mathrm{cur}}$, and $\ell_{i}$ using $\mathbf{S}_{\mathrm{cur}}$; update the scaling of $\mathbf{O}_i$ accordingly.
      \FOR{$1 \le j < T_c - 1$}
        \STATE Wait for $\mathbf{K}_j$ to be transferred into shared memory.
        \label{code:mainloop_start}
        \STATE Launch WGMMA to produce $\mathbf{S}_{\mathrm{next}} = \mathbf{Q}_i \mathbf{K}_{j}^T$, then commit the computation but do not wait for completion. \label{code:first_wgmma}
        \STATE Wait until $\mathbf{V}_{j-1}$ is resident in shared memory.
        \STATE Update $\mathbf{O}_{i}$ by adding $\tilde{\mathbf{P}}_{\mathrm{cur}} \mathbf{V}_{j-1}$ via WGMMA. Commit, do not wait. \label{code:second_wgmma}
        \STATE Wait for the WGMMA result of $\mathbf{Q}_i \mathbf{K}_{j}^T$.
        \STATE Recompute $m_{i}$, $\tilde{\mathbf{P}}_{\mathrm{next}}$, and $\ell_{i}$ from $\mathbf{S}_{\mathrm{next}}$. \label{code:softmax}
        \STATE Wait for $\tilde{\mathbf{P}}_{\mathrm{cur}} \mathbf{V}_{j-1}$ WGMMA, then re-apply scaling to $\mathbf{O}_i$
        \STATE Release the $(j\,\%\,s)$th and $(j-1\,\%\,s)$th sections of the buffers tied to $\mathbf{K}$ and $\mathbf{V}$, respectively.
        \STATE Assign the contents of $\mathbf{S}_{\mathrm{next}}$ to $\mathbf{S}_{\mathrm{cur}}$.
        \label{code:mainloop_end}
      \ENDFOR \STATE Wait for $\mathbf{V}_{T_c - 1}$ to become available in shared memory.
      \STATE Employ WGMMA to compute the final output: $\mathbf{O}_{i} = \mathbf{O}_{i} + \tilde{\mathbf{P}}_{\mathrm{last}} \mathbf{V}_{T_c - 1}$, then commit and wait for all operations to complete.

\STATE Epilogue: Adjust $\mathbf{O}_{i}$ according to $m_{i}$. Determine $L_{i}$ utilizing both $m_{i}$ and $\ell_{i}$.
      Store $\mathbf{O}_{i}$ and $L_{i}$ in HBM, assigning them to the $i$-th segment of $\mathbf{O}$ and $L$ respectively.
    \end{algorithmic}
  \end{algorithm}

\cref{alg:flash3_wgmma} substitutes the consumer path presented in \cref{alg:flash3_wgmma_ws_only}, thereby enabling the full implementation of the \textsc{FlashAttention-3} algorithm at FP16 precision. Conceptually, the acronym WGMMA in this context is used to denote asynchronous GEMM operations. During the mainloop (spanning lines \ref{code:mainloop_start} through \ref{code:mainloop_end}), the algorithm schedules the second WGMMA computation of the current iteration $j$ (as specified in line \ref{code:second_wgmma}) so that it runs concurrently with the softmax computations of the subsequent iteration $j+1$ (line \ref{code:softmax}).

Although the pipelined architecture depicted previously provides potential performance improvements in theory, several practical considerations arise:
\paragraph{Compiler reordering} The pseudocode assumes an ideal sequence of operations; however, the NVCC compiler frequently reorders instructions for optimization purposes. Such rearrangement can interfere with the intended alternation of WGMMA and non-WGMMA operations within the pipeline, potentially undermining performance advantages or causing unanticipated effects. Examination of the SASS code indicates that the compiler does, in fact, produce overlapped instructions as intended (see Section~\ref{sec:2-stage-sass}).
\paragraph{Register pressure} Achieving peak efficiency requires keeping register spilling to a minimum. Nevertheless, implementing the 2-stage pipeline entails the use of additional registers to hold temporary values and retain inter-stage context. In particular, a supplementary $\mathbf{S}_{\mathrm{next}}$ variable must reside in registers, contributing an extra overhead of $B_r \times B_c \times \text{sizeof}(\text{float})$ in register allocation per threadblock. This heightened register consumption can compete with other optimizations, such as increasing block dimensions, which similarly demand considerable register resources. Ultimately, optimization decisions regarding block size and pipeline configuration should be informed by empirical profiling.
\paragraph{3-stage pipelining} Building upon the previously discussed 2-stage pipeline, a 3-stage extension is introduced that further overlaps the second WGMMA with the softmax operation. This approach holds promise for still greater utilization of Tensor Cores but simultaneously demands even more register resources to accommodate the additional pipeline stage. This requirement complicates the balance between selecting suitable tile sizes and determining optimal pipeline depth. The specifics of the 3-stage algorithm and a discussion of its empirical evaluation are detailed in ~\cref{sec:3-stage}.

\subsection{Low-precision with FP8}
\label{sec:algofp8}

\textbf{Efficiency: layout transformations.} Executing the forward computation of \textsc{FlashAttention-3} using FP8 precision introduces unique difficulties with respect to layout compatibility that do not arise when operating in FP16.

To begin, it is important to observe that the input tensors $\mathbf{Q}$, $\mathbf{K}$, and $\mathbf{V}$ are commonly stored contiguously along the head dimension. However, for the second GEMM on FP8 WGMMA—which enforces a k-major constraint—it is necessary that $\mathbf{V}$, specifically the tiles transferred into SMEM, be contiguous along the sequence length dimension. Since the TMA load process itself does not alter which dimension is contiguous, there are two possible strategies: (1) transpose $\mathbf{V}$ in GMEM in advance, or (2) perform an in-kernel transpose on the loaded tiles of $\mathbf{V}$ in SMEM. For the first approach, this can be accomplished either by (1a) fusing the transposition step with the epilogue of an earlier computation stage, such as the rotary embedding, or by (1b) invoking a dedicated pre-processing transpose kernel\footnote{A high-performance transpose kernel can achieve throughput close to the theoretical memory bandwidth of the hardware~\citep{colfax_cutlass_transpose_2024}.} that swaps the head and sequence length dimensions. However, the fused strategy in (1a) tends to be challenging to incorporate within a general-purpose library, while the standalone transpose in (1b) is generally inefficient for inference workloads, which are often constrained by memory bandwidth.

For FP8 \textsc{FlashAttention-3}, we instead choose approach (2). During the in-kernel transpose, we leverage the LDSM (\verb|ldmatrix|) and STSM (\verb|stmatrix|) instructions, which permit a warp of threads to collectively transfer data between SMEM and RMEM in 128-byte chunks. \footnote{While the PTX documentation characterizes LDSM/STSM as copying $8 \times 8$ matrices composed of 16-bit elements \cite[\S 9.7.13.4.15-16]{ptx}, it is possible to employ these instructions with FP8 by packing pairs of 8-bit elements. However, the transpose variants of LDSM/STSM lack the ability to separate packed 8-bit values, making additional register shuffles necessary between the LDSM and STSM instructions to accomplish the required tile-level transpose; the full implementation details are omitted here.} These LDSM/STSM primitives are efficient in their use of registers—enabling execution within the producer warpgroup—and support layout transpositions during memory transfers. Additionally, once the first iteration concludes, we can overlap the transposition of the subsequent $\mathbf{V}$ tile with the two WGMMAs that process the previous $\mathbf{V}$ tile and the active $\mathbf{K}$ tile.

Second, it should be noted that the FP32 accumulator associated with an FP8 WGMMA employs a memory layout distinct from that utilized for operand A when stored in registers, which stands in contrast to the behavior observed with FP16. Illustrative portions of these two memory layouts are provided in \cref{fig:rmem_accum} and \cref{fig:rmem_operand}, where entries assigned to individual threads are listed in the order corresponding to their storage in registers. By leveraging byte permutation instructions, the accumulator output of the initial WGMMA can be reorganized into a layout compatible both with the requirements of a subsequent WGMMA and the structure of the $\mathbf{V}$ tile, as constructed by the in-kernel transpose operation. In particular, taking \cref{fig:rmem_accum} as reference, we implement a sequential rearrangement yielding
$$\{ \verb|d0 d1 d4 d5 d2 d3 d6 d7| \},$$
and then repeat this pattern for every consecutive group of 8 bytes. Interpreted in terms of the logical arrangement of the $\mathbf{P}$ tile, this action effectively permutes columns (e.g., columns $0189$ are relocated to become the leading four columns). To ensure the WGMMA delivers the correct result for the output tile, the in-kernel transpose can accordingly be configured to apply a corresponding row reordering to the $\mathbf{V}$ tile.\footnote{Notably, the ability to coordinate the in-kernel transpose in this manner obviates the need for explicit shuffle instructions to transfer register values between threads—a requirement highlighted previously in~\citep{colfax_fp8_flashattention_2024}.}

\textbf{Accuracy: block quantization and incoherent processing.} The FP8 (e4m3) data format allocates only 3 bits for the mantissa and 4 bits for the exponent, which leads to greater numerical inaccuracies compared to FP16 or BF16 precision. Additionally, large-scale models often exhibit extreme outlier values~\citep{dettmers2208llm,
  sun2024massive} that can overshadow the majority of other elements, presenting challenges for effective quantization. Common practice involves per-tensor scaling~\citep{micikevicius2022fp8}, which assigns a single scaling factor to each tensor (for example, separate scalars for $\mathbf{Q}$, $\mathbf{K}$, and $\mathbf{V}$).
To address the heightened quantization error in FP8 attention, we incorporate two specific methods:

\begin{enumerate}

\item \textbf{Block quantization}: In this approach, we assign a single scaling factor to each block; specifically, for each of $\mathbf{Q}$, $\mathbf{K}$, and $\mathbf{V}$, the tensors are partitioned into blocks of dimensions $B_r \times d$ or $B_c \times d$, and each block is quantized independently. This quantization process can be combined with certain pre-attention operations (such as rotary embeddings) without incurring extra computational overhead, since such operations are typically limited by memory bandwidth. Given that the \textsc{FlashAttention-3} algorithm is inherently block-structured, it is possible to adjust the scaling of each block of $\mathbf{S}$ accordingly for block quantization, effectively introducing no additional computational expense.

\item \textbf{Incoherent processing}: To mitigate the impact of outlier values, both $\mathbf{Q}$ and $\mathbf{K}$ are premultiplied by a random orthogonal matrix $\mathbf{M}$ prior to quantization to FP8 precision. The orthogonality of $\mathbf{M}$ ensures that $\mathbf{M} \mathbf{M}^\top = I$ and thus $(\mathbf{Q} \mathbf{M})(\mathbf{K} \mathbf{M})^\top = \mathbf{Q}\mathbf{K}^\top$, indicating that this transformation does not affect the computed attention. This technique redistributes outliers by making each entry of $\mathbf{Q}\mathbf{M}$ or $\mathbf{K}\mathbf{M}$ a random linear combination of the respective inputs, thereby lowering quantization errors. Following the approach of \citet{chee2024quip} and \citet{tseng2024quip}, $\mathbf{M}$ is selected as the product of random diagonal sign matrices (with entries $\pm 1$) and a Hadamard matrix, which supports efficient multiplication in $O(d \log d)$ time rather than $O(d^2)$, and this step can similarly be fused with the rotary embedding without increasing computational cost.

We empirically demonstrate in \cref{sec:numerical_error} that these two strategies reduce numerical error by as much as 2.6$\times$.

\section{Empirical Validation}
\label{sec:exp}

Our implementation of \textsc{FlashAttention-3} leverages CUTLASS~\citep{Thakkar_CUTLASS_2023} primitives, specifically the WGMMA and TMA abstractions, for assessment of both performance and precision.

\begin{itemize}

\item \textbf{Benchmarking attention.} We evaluate the execution time of \textsc{FlashAttention-3} over a range of sequence lengths, making comparisons against the default PyTorch implementation, \textsc{FlashAttention-2}, the Triton-based \textsc{FlashAttention-2} which leverages H100-optimized instructions, as well as a cuDNN-based, vendor-tuned version of \textsc{FlashAttention-2} tailored for H100 GPUs. Our findings indicate that \textsc{FlashAttention-3} delivers up to a 2.0$\times$ increase in speed relative to \textsc{FlashAttention-2}, and is 1.5$\times$ faster than the Triton variant. Additionally, \textsc{FlashAttention-3} achieves peak throughput of 740 TFLOPs/s, corresponding to 75\% of the H100 GPU's theoretical peak performance.

\item \textbf{Ablation study.} We demonstrate through ablation experiments that our optimizations, specifically warp-specialization and the pipelining of GEMM and softmax operations, are major contributors to the performance gains observed in \textsc{FlashAttention-3}.

\item \textbf{Accuracy of FP8 attention.} We show that utilizing block quantization and incoherent computation results in a 2.6$\times$ reduction in numerical errors for FP8 \textsc{FlashAttention-3}.

\subsection{Benchmarking Attention}
\label{subsec:benchmark_attn}

We evaluate the execution time of various attention algorithms on an H100 80GB SXM5 GPU under multiple configurations, considering both the presence and absence of causal masks, as well as head dimensions of 64 and 128, using FP16 data. Our findings, detailed in~\cref{fig:benchmark_attn_fwd} and~\cref{fig:benchmark_attn_bwd}, indicate that \textsc{FlashAttention-3} achieves a forward pass speedup of approximately 1.5-2.0$\times$ and a backward pass speedup of roughly 1.5-1.75$\times$ over \textsc{FlashAttention-2}. When compared to conventional attention mechanisms, \textsc{FlashAttention-3} demonstrates improvements of up to 3-16$\times$ in runtime. Notably, for sequences of medium to substantial length (1k tokens or more), \textsc{FlashAttention-3} even outperforms a proprietary library (cuDNN, closed source) optimized for the H100 architecture.

\paragraph{Benchmark settings:} We vary the sequence length as 512, 1k, ..., 16k, and set batch size so that the total number of tokens is 16k. We set the hidden dimension to 2048, and head dimension to be either 64, 128, or 256 (i.e., 32 heads, 16 heads, or 8 heads). To calculate the FLOPs of the forward pass, we use:

\begin{equation*}
  4 \cdot \text{seqlen}^2 \cdot \text{head dimension} \cdot \text{number of heads}.
\end{equation*}

When applying causal masking, this value is halved to reflect that roughly 50% of the computations are performed. For the backward pass, the FLOPs are computed as 2.5 times those of the forward pass, because, owing to recomputation, there are 2 matrix multiplications during the forward pass and 5 during the backward pass.

Additionally, we assess the forward pass runtime using FP8 under comparable conditions. The outcomes for headdim 256 are presented in ~\cref{fig:benchmark_attn_fp8}, with a comprehensive set of results provided in \cref{sec:benchmark_attn_fp8_full}.

\subsection{Ablation Study: 2-Stage Pipelining Experiments}
\label{sec:2-stage-pipelining-experiments}

We evaluate the impact of removing both the 2-stage WGMMA-softmax pipelining and warp-specialization components in non-causal FP16 \textsc{FlashAttention-3}, using constant parameters $\{\text{batch}, \text{seqlen}, \text{nheads}, \text{hdim}\} = \{ 4, 8448, 16, 128\}$. As demonstrated in~\cref{table:ablation_pipelining}, these algorithmic advancements—specifically, implementing asynchrony with warp-specialization and enabling overlap between GEMM and softmax—substantially enhance performance, increasing throughput from 570 to 661 TFLOPs.

\subsection{Numerical Error Validation}
\label{sec:numerical_error}

Given recent attention to the issue of numerical error in \textsc{FlashAttention}~\citep{golden2024flash}, we carry out a comparison among \textsc{FlashAttention-2}, \textsc{FlashAttention-3}, and a baseline attention implementation using a high-precision FP64 reference. To reflect the presence of outlier features and activations observed in LLMs~\citep{dettmers2208llm, sun2024massive}, we sample the elements of $\mathbf{Q}, \mathbf{K}, \mathbf{V}$ according to the following distribution:

\begin{equation*}
  \mathcal{N}(0, 1) + \mathcal{N}(0, 100) \cdot \mathrm{Bernoulli}(0.001).
\end{equation*}

In this setup, each matrix element follows a normal distribution with a mean of zero and a standard deviation of one, except for 0.1\% of the elements, to which we add an independent normal variable with a standard deviation of ten. We assess the root mean squared error (RMSE) as reported in~\cref{table:numerical_error}. For FP16, both \textsc{FlashAttention-2} and \textsc{FlashAttention-3} yield an RMSE that is 1.7$\times$ lower relative to the conventional implementation, owing to the retention of intermediate values (softmax outputs) in FP32 precision. The standard FP8-based attention baseline utilizes per-tensor scaling, with matrix multiplication accumulating results in FP32 and intermediate softmax computations in FP16. Due to block quantization and the method of incoherent processing, \textsc{FlashAttention-3} operating in FP8 demonstrates RMSE that is 2.6$\times$ smaller compared to this FP8 baseline.

\section{Dicussion, Limitations, Conclusion}
\label{sec:discussion}

Through the development of \textsc{FlashAttention-3}, we have shown that leveraging modern programming methodologies alongside advancements in hardware—specifically, asynchronous execution and lower-precision arithmetic—can substantially enhance both the speed and precision of attention mechanisms. Our method achieves a 1.5-2.0$\times$ acceleration over \textsc{FlashAttention-2} and decreases FP8 numerical errors by a factor of 2.6 relative to traditional per-tensor quantization. Notwithstanding these gains, there are several avenues for future improvement, including optimizing the system specifically for LLM inference, incorporating a persistent kernel approach within the FP8 kernel,\footnote{Within our experimental setup, the FP16 version of \textsc{FlashAttention-3} employs a persistent kernel alongside an effective load-balancing scheme, whereas the FP8 counterpart lacks this optimization. This contributes to the observed inferior performance of FP8 \textsc{FlashAttention-3} on shorter sequence lengths and causal masking when contrasted with FP8 cuDNN kernels.} and further examining how reduced-precision attention operates in the context of large-scale model training. While our empirical evaluation has concentrated on Hopper GPUs, we anticipate the strategies introduced here to be transferable to other types of hardware accelerators. Ultimately, we envision that a more rapid and precise attention implementation may pave the way for advances in tasks requiring extended context.

\end{document}